\documentclass[a4paper]{article}
\usepackage{bm}
\usepackage{graphicx}
\usepackage[cm]{fullpage}

\author{C. Thieulot}
\title{GHOST}

\begin{document}

\maketitle

\begin{center}
\includegraphics[width=4cm]{ghost}
\end{center}

\tableofcontents

Please contact me if you need a new feature or if you wish to contribute to GHOST at 
c.thieulot@uu.nl

%---------------------------------
\section{Introduction}

The open source code library {\sc ghost} allows 
three different types of hollow sphere meshes to be built
, i.e. meshes bounded by two concentric spheres:

\begin{itemize}
\item 
The cubed sphere ('HS06'), composed of 6 blocks which 
are themselves subdivided into $N_b \times N_b$ quadrilateral shaped cells  \cite{sado72,roip96,heta03,busa13}.
%3 types: conformal \cite{rapm96,mcgr96}, obtained with elliptic solvers/spring dynamics \cite{puli07}, 
%gnomonic  \cite{sado72,mcgr96}
Four types of cubed spheres meshes have been proposed: the conformal, elliptic, gnomonic and spring types \cite{puli07}.
However only gnomonic meshes are considered here: these are obtained by inscribing a cube within a sphere and expanding to the surface
of the sphere.
The cubed sphere has recently been used in large-scale mantle convection simulation in conjunction with 
Adaptive Mesh Refinement \cite{algs12,busa13}.  

\item 
The CitcomS mesh ('HS12') composed of 12 blocks also subdivided 
into $N_b \times N_b$ quadrilateral shaped cells
\cite{zhzm00,sthh06,zhmt08,arfw14}.
Note that ASPECT \cite{krhb12,hedg17}, a relatively new code aimed at superseeding CitcomS can generate and use 
this type of mesh \cite{thie17} but is not limited to it.

\item The icosahedral mesh ('HS20') composed of 20 triangular blocks \cite{bafr85,baum85} subdivided into triangles, which is 
used in the TERRA code \cite{burb96,burb97,burl98,dadb13}.
\end{itemize}

Given the regularity and symmetry of these meshes determining the location of the 
mesh nodes in space is a relatively straightforward task. Building the mesh connectivity in an 
efficient manner is where the difficulty lies.

The approach to building all three meshes is identical:
\begin{enumerate}
\item A reference square or triangle is populated with cells, as shown in Fig. (\ref{basics}) 
parametrised by a level $l$: the square is subdivided into $l\times l$ quadrilaterals while 
the triangle is subdivided into $l^2$ triangles.

\item This reference square or triangle is then replicated {\sl nblock} times (6, 12 or 20) and mapped
onto a portion of a unit sphere. The blocks are such that their union covers a full sphere
but they cannot overlap except at the edges, see Fig. (\ref{3shell}).

\item All block meshes are then merged together to generate a shell mesh. 
\item Shell meshes are replicated {\sl nlayer+1} times outwards with increasing radii. 
\item The {\sl nlayer} shells are then merged together to form a hollow sphere mesh, as shown in 
Fig. (\ref{balls}).
\end{enumerate}

In Table (\ref{tabv}) the number of nodes and cells for a variety of resolutions 
for all three mesh types is reported. Looking at the CitcomS literature of the past 20 years, we find that 
the mesh data presented in this table cover the various resolutions used, e.g.
$12\times48^3$ \cite{mczh04,arfw14}, $12\times64^3$ \cite{budt14}
$12\times96^3$ \cite{bumb10}, $12\times128^3$ \cite{beck06,wele16,welm16}.
Note that in the case of the HS06 and HS12 meshes the mesh nodes are mapped out to the 6 or 12 blocks 
following either an equidistant or equiangle approach as shown in Fig. (\ref{equis2}) (see \cite{puli07}
for details on both approaches). 

\begin{figure}[t]
\centering
\includegraphics[width=12cm]{images/basic_shapes/basics}
\caption{Reference square and triangles meshes at level 5 \label{basics}}
\end{figure}

\begin{figure}[t]
\centering
\includegraphics[width=.3\linewidth]{images/3HS/shell_HS06}
\includegraphics[width=.3\linewidth]{images/3HS/shell_HS12}
\includegraphics[width=.3\linewidth]{images/3HS/shell_HS20}
\caption{From left to right: HS06, HS12 and HS20 shells coloured by block number. \label{3shell}}
\end{figure}

\begin{figure}[t]
\centering
\includegraphics[width=12cm]{images/3HS/3HS}
\caption{a) HS06 mesh composed of 6 blocks containing each $6^3$ cells; b) HS12 mesh composed of 12 blocks containing each 
$6^3$ cells; e) HS20 mesh composed of 20 blocks containing each $6^3$ cells. \label{balls}}
\end{figure}











%---------------------------------
\section{Conventions \& notations}

On the following figure are represented the three cartesian axis, 
a point and its spherical coordinates $r,\theta,\phi$:
\begin{center}
\includegraphics[width=5cm]{images/sphcoord}\\
{\small Spherical coordinates as commonly used in physics:\\ polar angle $\theta$, and azimuthal angle $\phi$.} 
\end{center}
In this case $\theta\in[0:\pi]$ and $\phi\in]-\pi:\pi]$ and we have the following relationships:
\begin{eqnarray}
r &=& \sqrt{x^2+y^2+z^2} \\
\theta &=& acos (z/r) \\
\phi &=& atan (y/x) 
\end{eqnarray}
The inverse tangent used to compute $\phi$ must be suitably defined, taking into account the correct quadrant of $(x,y)$,
which is why the atan2 intrinsic function is used.    

In geography one uses latitude and longitude, represented hereunder:
\begin{center}
\includegraphics[width=10cm]{images/map.jpg}
\end{center}
\begin{itemize}
\item Latitude  $\in[-90:90]$,   or $\in[-\pi/2:\pi/2]$ 
\item Longitude $\in]-180:180]$, or $\in]-\pi:\pi]$ 
\end{itemize}


%---------------------------------
\section{Running the program}

The code can be compiled by typing {\sl make}. The default compiler is {\sl gfortran}
but it can be changed by editing the {\sl Makefile}. 
Running {\sl make} generates the {\sl ghost} executable which can then be run as follows
in a terminal:

\begin{verbatim}
> ./ghost
\end{verbatim}

If you wish to know which options are available, simply do 
\begin{verbatim}
> ./ghost --help
\end{verbatim}

Options are as follows:
\begin{itemize}
\item {\tt -level X} where X is an integer number strictly larger than 2.
\item {\tt -nlayer X} where X is an integer number strictly larger than 2.
\item {\tt -mtype X} where X is equal to 6,12,20
\item {\tt -equiangular X} where X is 0 (false) or 1 (true) 
\end{itemize}

For instance, if you wish to generate a HS06 mesh of block level 16
with 20 layers in the radial direction with equiangular spacing:
\begin{verbatim}
./ghost -level 16 -mtype 6 -nlayer 20 -equiangular 1
\end{verbatim}

Note that a clean option exists to remove all compilation files and output files:
\begin{verbatim}
> make clean 
\end{verbatim}

%Data presented in the paper were obtained by running the shell script:
%\begin{verbatim}
%> ./script_paper 
%\end{verbatim}

\newpage

You can then open the file {\sl hollow\_sphere.vtu} in Paraview:


\includegraphics[width=6cm]{images/tas/r}
\includegraphics[width=6cm]{images/tas/theta}
\includegraphics[width=6cm]{images/tas/phi}


\includegraphics[width=6cm]{images/tas/inner}
\includegraphics[width=6cm]{images/tas/outer}
\includegraphics[width=6cm]{images/tas/blocknb}

\includegraphics[width=8cm]{images/tas/volume}
\includegraphics[width=8cm]{images/tas/normals}

\newpage


\section{Subroutines}
 \subsection{\tt allocate\_block\_memory}
 This subroutine assigns to every block {\sl i} a number of node 
 points {\sl block(i)\%np} and a number of cells {\sl block(i)\%ncell}  
 The number of vertices per cell is also stored in {\sl block(i)\%nv}.
 It then loops over all blocks and allocates the required arrays 
 to store the mesh nodes position in Cartesian and spherical coordinates
 and the connectivity.  
 \subsection{\tt block\_node\_layout}
 This subroutine generates the position of the mesh node points for each 
 block as well as their connectivity (i.e. the list of nodes making 
 each element. 
 \begin{center}
 \includegraphics[width=8cm]{images/basics}\\
 {\small a) block layout for mtype=3,4; b) block layout for mtype=5}
 \end{center}
 In the case that the block is a quadrilateral, it is assumed that 
 it is made of $l \times l$ quadrilateral cells. In the case that the block is a 
 triangle, it is made of $l^2$ triangular cells.
 \subsection{\tt build\_hollow\_sphere}
 This subroutine first makes a copy of the current shell into {\sl shell\_temp}
 It then computes the number of mesh nodes and cells for the hollow sphere
 and allocate all arrays accordingly.
 It then loops over layers, places the temporary shell at the required 
 radius, and appends it to the current hollow sphere mesh from the inside out.
 \subsection{\tt check\_parameters}
 \subsection{\tt }
 \subsection{\tt hexahedron\_volume}
 This function computes the volume of any hexahedron, 
 following Eq.(12) of \cite{gran97}.
 \subsection{\tt compute\_gravity\_at\_point}
 This subroutine receives as argument the coordinates {\tt x,y,z} of a point and 
 returns the gravity vector components {\tt gx,gy,gz} at this location 
 by means of 
 \[
 {\bm g}({\bm r}') = \int_\Omega {\cal G} \frac{\rho({\bm r})-\rho_0}{|{\bm r}'-{\bm r}|^3} ({\bm r}'-{\bm r}) d\Omega
 \] 
 where $\rho_0=${\tt refdensgrav}.
 The gravitational potential is computed as follows:
 \[
 U({\bm r}') = -\int_\Omega {\cal G} \frac{\rho({\bm r})-\rho_0}{|{\bm r}'-{\bm r}|}  d\Omega
 \]
 \subsection{\tt compute\_gravity\_on\_line}
 This subroutine computes on a line parametrised by $0\leq r \leq 2R_2$,
 $\theta=13^\circ$ and $\phi=17^\circ$. This line is discretised 
 over npts points. At each point is the subroutine compute\_gravity\_at\_point called.
 \subsection{\tt compute\_hollow\_sphere\_volume}
 \subsection{\tt compute\_normals\_hollow\_sphere}
 \subsection{\tt compute\_r\_theta\_phi\_hollow\_sphere}
 This subroutine computes $r$, $\theta$ and $\phi$ from $x,y,z$ for all 
 nodes of the hollow\_sphere mesh with
 \[
 r=\sqrt{x^2+y^2+z^2}
 \]
 \[
 \theta=\cos^{-1}(z/r)
 \]
 \[
 \phi=\tan^{-1}_2 (y/x)
 \]
 \subsection{\tt compute\_r\_theta\_phi\_shell}
 This subroutine computes $r$, $\theta$ and $\phi$ from $x,y,z$ for all 
 nodes of the shell mesh with
 \[
 r=\sqrt{x^2+y^2+z^2}
 \]
 \[
 \theta=\cos^{-1}(z/r)
 \]
 \[
 \phi=\tan^{-1}_2 (y/x)
 \]
 \subsection{\tt compute\_shell\_area}
 \subsection{\tt compute\_shell\_distributions}
 \subsection{\tt compute\_triangle\_area}
 {\tt https://en.wikipedia.org/wiki/Heron\%27s\_formula }
 Heron's formula states that the area of a triangle whose sides have lengths a, b, and c is
 \[
 A=\sqrt{s(s-a)(s-b)(s-c)}
 \]
 where s is the semiperimeter of the triangle; that is,
 \[
 s=\frac{a+b+c}{2}
 \]
 \subsection{\tt generate\_quadrature\_points}
 \begin{center}
 \includegraphics[width=6cm]{images/elements}
 \end{center}
 For hexahedral elements the linear shape functions are:
 \begin{eqnarray}
 N_1(r,s,t) &=& (1-r)(1-s)(1-t)/8 \nonumber\\ 
 N_2(r,s,t) &=& (1+r)(1-s)(1-t)/8 \nonumber\\ 
 N_3(r,s,t) &=& (1+r)(1+s)(1-t)/8 \nonumber\\ 
 N_4(r,s,t) &=& (1-r)(1+s)(1-t)/8 \nonumber\\ 
 N_5(r,s,t) &=& (1-r)(1-s)(1+t)/8 \nonumber\\ 
 N_6(r,s,t) &=& (1+r)(1-s)(1+t)/8 \nonumber\\ 
 N_7(r,s,t) &=& (1+r)(1+s)(1+t)/8 \nonumber\\ 
 N_8(r,s,t) &=& (1-r)(1+s)(1+t)/8 \nonumber
 \end{eqnarray}
 For the triangular prism elements the linear shape functions are:
 \begin{eqnarray}
 N_1(r,s,t) &=& (1-r-s)(1-t)/2 \nonumber\\ 
 N_2(r,s,t) &=& r(1-t)/2 \nonumber\\ 
 N_3(r,s,t) &=& s(1-t)/2 \nonumber\\ 
 N_4(r,s,t) &=& (1-r-s)(1+t)/2 \nonumber\\ 
 N_5(r,s,t) &=& r(1+t)/2 \nonumber\\ 
 N_6(r,s,t) &=& s(1+t)/2 \nonumber
 \end{eqnarray}
 \subsection{\tt map\_blocks}
 \subsubsection{HS06 - cubed sphere}
 See section 3.1 of \cite{puli07}
 \subsubsection{HS12 - CitcomS mesh}
 \subsubsection{HS20 - Icosahedron}
 Let $t=(1+\sqrt{5})/2$, then the vertices of the icosahedron are 
 \[
 {\bm x}_0=\frac{1}{\sqrt{1+t^2}}(t,1,0) \quad\quad {\bm x}_6=\frac{1}{\sqrt{1+t^2}}(-1,0,t)
 \]
 \[
 {\bm x}_1=\frac{1}{\sqrt{1+t^2}}(-t,1,0) \quad\quad {\bm x}_7=\frac{1}{\sqrt{1+t^2}}(-1,0,-t)
 \]
 \[
 {\bm x}_2=\frac{1}{\sqrt{1+t^2}}(t,-1,0) \quad\quad {\bm x}_8=\frac{1}{\sqrt{1+t^2}}(0,t,1)
 \]
 \[
 {\bm x}_3=\frac{1}{\sqrt{1+t^2}}(-t,-1,0) \quad\quad {\bm x}_9=\frac{1}{\sqrt{1+t^2}}(0,-t,1)
 \]
 \[
 {\bm x}_4=\frac{1}{\sqrt{1+t^2}}(1,0,t) \quad\quad {\bm x}_{10}=\frac{1}{\sqrt{1+t^2}}(0,t,-1)
 \]
 \[
 {\bm x}_5=\frac{1}{\sqrt{1+t^2}}(1,0,-t) \quad\quad {\bm x}_{11}=\frac{1}{\sqrt{1+t^2}}(0,-t,-1)
 \]
 It is obvious that all points lie on a sphere of unit radius.

 The triangles are as follows:
 \[  T_0=[0,8,4]  \quad\quad T_{10}=[2,9,11]   \]
 \[  T_1=[0,5,10] \quad\quad T_{11}=[3,11,9]   \]
 \[  T_2=[2,4,9]  \quad\quad T_{12}=[4,2,0]   \]
 \[  T_3=[2,11,5] \quad\quad T_{13}=[5,0,2]   \]
 \[  T_4=[1,6,8]  \quad\quad T_{14}=[6,1,3]   \]
 \[  T_5=[1,10,7] \quad\quad T_{15}=[7,3,1]   \]
 \[  T_6=[3,9,6]  \quad\quad T_{16}=[8,6,4]   \]
 \[  T_7=[3,7,11] \quad\quad T_{17}=[9,4,6]   \]
 \[  T_8=[0,10,8] \quad\quad T_{18}=[10,5,7]   \]
 \[  T_9=[1,8,10] \quad\quad T_{19}=[11,7,5]   \]
 Graphically, the mesh has the following structure:
 \begin{center}
 \includegraphics[width=12cm]{images/flat_icosahedron}
 \end{center}
 \subsection{\tt merge\_blocks}
 \subsection{\tt output\_blocks\_vtu}
 \subsection{\tt output\_hollow\_sphere\_vtu}
 \subsection{\tt output\_shell\_ascii}
 \subsection{\tt output\_shell\_lat\_lon}
 This subroutine exports $\theta$ and $\phi$ values of all points on the 
 shell in the file {\sl shell\_theta\_phi.dat} in {\sl OUT}.
 \subsection{\tt output\_shell\_vtu}
 If {\sl generate\_vtu\_output} is true 
 a vtu file of the shell is created in {\sl OUT}.
 \subsection{\tt project\_on\_sphere}
 This subroutine gets a radius $r$ value as argument and the current coordinates
 of a point. It returns the new coordinates of the point on a sphere 
 of radius $r$ with the same $\theta,\phi$.
 \subsection{\tt read\_command\_line\_options}
 This subroutine reads in the command line arguments
 \subsection{\tt read\_s40rts}
 \subsection{\tt }
 \subsection{\tt write\_positions}
 \subsection{\tt write\_field\_dp}
 \subsection{\tt write\_field\_int}
 \subsection{\tt write\_icon}
 \subsection{\tt write\_offsets}
 \subsection{\tt write\_types}


\bibliographystyle{plain}
\bibliography{../../../../../biblio_geosciences2}

\end{document}
